\section{Functional Requirements}

The functional capabilities of VisionLab are categorized across the platform's two core operational modes (Mode A: Image Analysis and Mode B: Video Analysis) alongside core platform export and architectural utilities:

\subsubsection{Mode A: Image Analysis}
\begin{itemize}
    \item \textbf{FR-1.1 (Low-Level Filtering \& Spatial Convolutions):} The system shall support image smoothing and enhancement operations including Gaussian Blur (kernel sizes $3\times3$ to $31\times31$), Median Filtering (window sizes $3\times3$ to $21\times21$), and Laplacian high-pass sharpening via discrete spatial convolution kernels:
    \begin{equation}
        K_{\text{Laplacian}} = \begin{bmatrix} 0 & -1 & 0 \\ -1 & 5 & -1 \\ 0 & -1 & 0 \end{bmatrix}
    \end{equation}
    \item \textbf{FR-1.2 (Intensity \& Histogram Operations):} The system shall generate pixel intensity distribution histograms across grayscale and individual RGB color channels, and perform global Histogram Equalization to enhance low-contrast inputs.
    \item \textbf{FR-1.3 (Thresholding Operations):} The system shall implement global and local adaptive thresholding, including Otsu's optimal intra-class variance minimization, Adaptive Mean thresholding, and Adaptive Gaussian thresholding.
    \item \textbf{FR-1.4 (Mathematical Morphology):} The system shall provide morphological transformations including Dilation, Erosion, Opening, Closing, and Morphological Gradient using Rectangular, Elliptic, and Cross structuring elements.
    \item \textbf{FR-2.1 (Edge \& Gradient Detection):} The system shall extract image boundaries using Sobel horizontal/vertical first-order derivatives, Laplacian second-order derivatives, and the multi-stage Canny edge detector with tunable hysteresis thresholds.
    \item \textbf{FR-2.2 (Geometric Feature Detection):} The system shall compute the Probabilistic Hough Line Transform to detect collinear structures within gradient maps.
    \item \textbf{FR-2.3 (Corner \& Keypoint Detection):} The system shall compute Harris corner response maps based on the structure tensor $M$, and extract scale-invariant keypoints using Difference-of-Gaussians (DoG) SIFT feature extraction.
    \item \textbf{FR-2.4 (Dense Feature Descriptors):} The system shall extract Histogram of Oriented Gradients (HOG) feature vectors across normalized spatial cell grids.
    \item \textbf{FR-3.1 (Clustering-Based Segmentation):} The system shall segment images into $k$ color clusters using K-Means vector quantization in RGB/CIELAB space and Mean Shift density mode seeking.
    \item \textbf{FR-3.2 (Region \& Contour Segmentation):} The system shall support seeded 4-connected Region Growing and edge-based contour hierarchical polygon extraction.
    \item \textbf{FR-4.1 (Deep Learning Detection):} The system shall perform single-shot object detection using YOLOv8n, predicting bounding boxes, class labels (80 COCO classes), and confidence probabilities.
    \item \textbf{FR-4.2 (Interactive NMS Tuning):} The system shall enable live manipulation of the confidence threshold ($0.10$ to $0.90$) and Intersection-over-Union (IoU) Non-Maximum Suppression threshold ($0.10$ to $0.90$).
\end{itemize}

\subsubsection{Mode B: Video Analysis}
\begin{itemize}
    \item \textbf{FR-5.1 (Video Decoding \& Stream Telemetry):} The system shall ingest standard video codecs (MP4, AVI, MOV, MKV) and extract resolution, frame rate (FPS), frame count, duration, and file size.
    \item \textbf{FR-5.2 (Keyframe Temporal Sampling):} The system shall extract $N$ evenly-spaced keyframes across the timeline and render a horizontal visual summary strip with frame index markers.
    \item \textbf{FR-6.1 (Multi-Object Tracking):} The system shall integrate YOLOv8n with the ByteTrack algorithm, maintaining persistent tracking IDs across temporal frames via Kalman filter state prediction and Hungarian matching.
    \item \textbf{FR-6.2 (Temporal Decimation):} The system shall allow users to adjust the frame step parameter ($1$ to $30$) to optimize the trade-off between tracking accuracy and processing latency.
    \item \textbf{FR-6.3 (Tracking Telemetry):} The system shall display an annotated keyframe grid, category frequency distribution charts, and a detailed tracking log table.
    \item \textbf{FR-7.1 (Dense Optical Flow):} The system shall compute Gunnar Farneback dense optical flow fields and visualize motion vector orientation and velocity via an HSV color wheel mapping.
    \item \textbf{FR-7.2 (Sparse Feature Tracking):} The system shall compute Lucas-Kanade pyramidal sparse optical flow (KLT) on Shi-Tomasi corner points.
    \item \textbf{FR-7.3 (Adaptive Background Subtraction):} The system shall separate moving foreground objects from stationary scenes using Mixture of Gaussians (MOG2) and K-Nearest Neighbors (KNN).
    \item \textbf{FR-7.4 (Motion Vector Analytics):} The system shall plot an 8-cardinal polar area direction histogram and a temporal motion magnitude time-series graph.
\end{itemize}

\section{Non-functional Requirements}
\begin{itemize}
    \item \textbf{NFR-1 (Performance \& Execution Latency):} Classical image filtering and feature extraction operations on standard $1080$p images shall execute in under $800$ milliseconds on standard consumer CPUs. Video processing pipelines shall support frame decimation to sustain processing speeds exceeding $15$ frames per second on CPU.
    \item \textbf{NFR-2 (Usability \& Design Aesthetics):} The user interface shall employ a high-contrast dark glassmorphism design system featuring semi-transparent panels (`\texttt{rgba(15, 22, 40, 0.86)}`), backdrop blur filters ($20$px), crisp borders, and accessible typography (`\texttt{\#f8fafc}` primary, `\texttt{\#cbd5e1}` secondary) ensuring complete readability across ambient shader backgrounds.
    \item \textbf{NFR-3 (Modularity \& Architectural Separation):} The backend shall strictly adhere to a modular architecture separating FastAPI endpoint routers (`\texttt{routers/}`), business logic and models (`\texttt{services/}`), and mathematical utilities (`\texttt{cv\_utils.py}`).
    \item \textbf{NFR-4 (Portability \& Zero-Dependency Frontend):} The frontend shall operate entirely using native browser technologies (Vanilla HTML5, CSS3, ES6 JavaScript) without requiring Node.js, Webpack, or external package managers, enabling immediate hosting via any static web server.
    \item \textbf{NFR-5 (Robustness \& Fault Tolerance):} The API shall validate file MIME types, sanitize multi-part payload boundaries, handle malformed binary headers gracefully, and return informative HTTP $4$xx/$5$xx JSON error contracts without crashing the ASGI daemon.
    \item \textbf{NFR-6 (Interoperability \& Documentation):} All endpoints shall expose formal OpenAPI 3.0 schemas automatically rendered via interactive Swagger UI at `\texttt{/api/docs}`.
\end{itemize}
